\section{Related Work}
\label{sec:related}

\paragraph{Memory mechanisms in agents.}
External memory helps LLM agents reuse experience beyond the context window. Early methods store raw trajectories as examples~\citep{zhao2024expel, chhikara2025mem0}, while later work compresses experience into summaries or knowledge entries~\citep{fang2025memp, liu2025simplemem, ouyang2025reasoningbank, tang2025agentkb}. Recent studies further apply RL directly to agent knowledge structures: MemRL~\citep{zhang2026memrl} performs runtime RL on episodic memory, MemEvolve~\citep{zhang2025memevolve} meta-evolves memory systems, Mem-$\alpha$~\citep{wang2025mem_alpha} learns memory construction policies, and EvolveR~\citep{wu2025evolver} co-adapts the policy and memory bank. In contrast, \textsc{SkillGraph} represents experience as explicit skill abstractions with typed dependencies and evolves this structure jointly with the policy.

\paragraph{Graph structures for LLMs.}
Graph structures have been widely adopted in LLM systems: Graph-of-Thought~\citep{besta2024got} models reasoning steps as a directed graph to enable non-linear thought exploration, GraphRAG~\citep{edge2024graphrag} builds entity-relation graphs over corpora for structured retrieval, and \citet{huang2024graphplan} encode task decompositions as planning graphs for agent execution. \textsc{SkillGraph} applies graph structures to agent skill management, jointly evolving the graph topology and the policy through RL, enabling the skill graph to adapt continuously rather than remaining static after construction.

\paragraph{Agent skill evolution.}
Agentic skills can compact reusable strategies for subtasks. Voyager~\citep{wang2023voyager} accumulates executable code skills, and ExpeL~\citep{zhao2024expel} distills transferable strategic experience from trajectories. Most closely related, SkillRL~\citep{xia2025skillrl} co-evolves a hierarchical skill bank with the agent's policy through recursive RL. \textsc{SkillGraph} builds on this line by elevating the flat skill bank into a structured dependency graph, enabling typed relational modeling and topology evolution throughout training.
